\documentclass[10pt,a4paper]{article}

\usepackage[spanish,es-nodecimaldot]{babel}
\usepackage[utf8]{inputenc}
\usepackage[T1]{fontenc}
\usepackage{newtxtext,newtxmath}
\usepackage{geometry}
\usepackage{microtype}
\usepackage{amsmath,mathtools}
\usepackage{enumitem}
\usepackage{xcolor}
\usepackage{hyperref}
\usepackage{fancyhdr}
\usepackage{lastpage}
\usepackage[most]{tcolorbox}

\geometry{margin=2.5cm,includeheadfoot}
\setlength{\headheight}{15pt}
\setlength{\headsep}{0.28cm}
\setlength{\footskip}{0.62cm}
\setlength{\parindent}{0pt}
\setlength{\parskip}{0.28em}
\setlength{\abovedisplayskip}{0.5em}
\setlength{\belowdisplayskip}{0.5em}
\setlength{\abovedisplayshortskip}{0.25em}
\setlength{\belowdisplayshortskip}{0.25em}
\setlist[enumerate]{leftmargin=2.2em,itemsep=0.3em,topsep=0.25em}
\setlist[itemize]{leftmargin=1.5em,itemsep=0.2em,topsep=0.2em}

\definecolor{repblue}{gray}{0}
\definecolor{repteal}{gray}{0}
\definecolor{repgray}{gray}{0}
\definecolor{repaccent}{gray}{0}
\definecolor{repline}{gray}{0}
\definecolor{replight}{gray}{1}

\hypersetup{
  hidelinks,
  pdftitle={TRGF Padilla Tarea 3},
  pdfauthor={Juan Ignacio Padilla Barrientos}
}

\pagestyle{fancy}
\fancyhf{}
\fancyhead[L]{\textcolor{repblue}{\small\itshape TRGF Padilla}}
\fancyhead[R]{\textcolor{repgray}{\small Seminario UCR I 2026}}
\fancyfoot[C]{\textcolor{repgray}{\small \thepage/\pageref*{LastPage}}}
\renewcommand{\headrulewidth}{0.3pt}
\renewcommand{\footrulewidth}{0pt}

\newcommand{\CC}{\mathbb{C}}
\newcommand{\ZZ}{\mathbb{Z}}

\tcbset{
  enhanced,
  sharp corners,
  boxrule=0.45pt,
  arc=0pt,
  left=0.85em,
  right=0.85em,
  top=0.45em,
  bottom=0.4em,
  boxsep=0pt
}

\newcommand{\inlinehead}[2][repblue]{%
  \par\smallskip
  {\color{#1}\bfseries #2.}\hspace{0.45em}%
}

\newcommand{\blockhead}[2][repteal]{%
  \par\smallskip
  {\color{#1}\bfseries #2.}\par\nobreak\smallskip
}

\newcommand{\answer}{\inlinehead{Respuesta}}
\newcommand{\subhead}[1]{\blockhead{#1}}
\newcommand{\notehead}[1]{\inlinehead[repgray]{#1}}

\newcommand{\topbanner}{%
  \begin{tcolorbox}[
    colback=replight,
    colframe=repline,
    borderline west={1.2pt}{0pt}{repblue},
    borderline east={1.2pt}{0pt}{repteal},
    left=1.0em,
    right=1.0em,
    top=0.7em,
    bottom=0.55em
  ]
    {\Large\bfseries\color{repblue} TRGF Padilla \textemdash\ Tarea 3}\hfill
    {\color{repaccent}\small\itshape Soluciones}\\[-0.1em]
    {\normalsize\color{repteal} Seminario UCR I 2026}
  \end{tcolorbox}
}

\newcounter{probcnt}
\newcommand{\problem}[2]{%
  \stepcounter{probcnt}%
  \ifnum\value{probcnt}>1\newpage\fi
  \par\vspace{0.25em}
  {\large\bfseries\color{repblue} Problema #1}\par
  {\color{repline}\rule{\linewidth}{0.55pt}}\par
  \begin{tcolorbox}[
    breakable,
    colback=white,
    colframe=repline,
    borderline west={1.6pt}{0pt}{repteal},
    left=0.7em,
    right=0.7em,
    top=0.35em,
    bottom=0.3em,
    before skip=0.35em,
    after skip=0.45em
  ]
    #2
  \end{tcolorbox}
}

\newcommand{\worklines}[1]{%
  \par\vspace{0.35em}
  {\color{repline}\rule{\linewidth}{0.4pt}}\par
  \begingroup
  \count0=1
  \loop
    \ifnum\count0<#1
      \vspace{0.72cm}
      {\color{repline}\rule{\linewidth}{0.4pt}}\par
      \advance\count0 by 1
  \repeat
  \endgroup
  \vspace{0.35em}
}

\begin{document}

\color{black}
\topbanner
\small

\problem{1}{%
\textit{\textquestiondown Para cuáles $n$ se cumple que hay un homomorfismo sobreyectivo de grupos $S_{n+1} \to S_n$?}
}

\answer
Existe un homomorfismo sobreyectivo $S_{n+1}\to S_n$ si y solo si
$n\in\{1,2,3\}$.

\subhead{Hecho que usaremos}

\notehead{Primer teorema de isomorfismo}
si $\phi\colon G\to H$ es un homomorfismo sobreyectivo de grupos,
entonces $G/\ker\phi\cong H$. En particular,
$|\ker\phi|=|G|/|H|$.

\subhead{Lema 1 (Jordan)}
Si $m\ge 5$, entonces $A_m$ es simple.

\subhead{Caso $n=1$}
$S_2\to S_1=\{e\}$. El único homomorfismo (el trivial) es sobreyectivo.

\subhead{Caso $n=2$}
El homomorfismo signo $\operatorname{sgn}\colon S_3\to\{+1,-1\}$ es sobreyectivo.
Como $\{+1,-1\}\cong \ZZ/2\ZZ\cong S_2$, obtenemos el homomorfismo sobreyectivo deseado.

\subhead{Caso $n=3$}
Considere las tres particiones de $\{1,2,3,4\}$ en dos pares:
\[
  P_1=\bigl\{\{1,2\},\{3,4\}\bigr\},\quad
  P_2=\bigl\{\{1,3\},\{2,4\}\bigr\},\quad
  P_3=\bigl\{\{1,4\},\{2,3\}\bigr\}.
\]
Estas son \emph{todas} las particiones de $\{1,2,3,4\}$ en dos pares
(hay $\binom{4}{2}/2!=3$).
Cada $\sigma\in S_4$ actúa sobre una partición aplicando $\sigma$ a cada elemento;
como $\sigma$ es biyección, el resultado es otra partición en dos pares,
y como solo hay tres, $\sigma$ permuta $\{P_1,P_2,P_3\}$; además,
$(1\;2)$ induce la transposición $(P_2\;P_3)$ y $(1\;2\;3)$ induce
el $3$-ciclo $(P_1\;P_3\;P_2)$.

\textit{Que $\phi$ es un homomorfismo:} esto es inmediato, pues una acción
de grupo sobre un conjunto $X$ es precisamente un homomorfismo
$G\to S(X)$; aquí $X=\{P_1,P_2,P_3\}$ y $S(X)\cong S_3$.

\textit{Sobreyectividad:}
la imagen de $\phi$ contiene una transposición y un $3$-ciclo en $S_3$.
En efecto, $\phi((1\;2))=(P_2\;P_3)$ y $\phi((1\;2\;3))=(P_1\;P_3\;P_2)$.
Un subgrupo de $S_3$ que contiene una transposición y un $3$-ciclo es todo $S_3$
(pues $|S_3|=6$ y el subgrupo generado tiene orden divisible por $2$ y por $3$, así que es $6$).
Por tanto $\phi$ es sobreyectivo.

\subhead{Caso $n\ge 4$}
Supongamos que existe $\phi\colon S_{n+1}\to S_n$ sobreyectivo.
Entonces $K:=\ker\phi$ es un subgrupo normal de $S_{n+1}$ de orden
\[
  |K|=\frac{|S_{n+1}|}{|S_n|}=\frac{(n+1)!}{n!}=n+1.
\]
Para $n\ge 4$ se tiene $n+1\ge 5$. Afirmamos que los únicos subgrupos
normales de $S_{n+1}$ son $\{e\}$, $A_{n+1}$ y $S_{n+1}$: en efecto,
si $N\trianglelefteq S_{n+1}$, entonces $N\cap A_{n+1}\trianglelefteq A_{n+1}$,
 y por el Lema 1 se tiene $N\cap A_{n+1}=\{e\}$ o $N\cap A_{n+1}=A_{n+1}$. Si $N\cap A_{n+1}=A_{n+1}$,
 entonces o bien $N=A_{n+1}$ o bien $N$ contiene una permutación impar, y en este último caso $N=S_{n+1}$.
 Si $N\cap A_{n+1}=\{e\}$ y $N$ contuviera una permutación impar $\sigma$, entonces $\sigma^2\in A_{n+1}\cap N=\{e\}$.
 Como $N$ es normal, todo conjugado $\tau$ de $\sigma$ pertenece a $N$;
 además $\sigma\tau$ es par (producto de dos impares), así que
 $\sigma\tau\in N\cap A_{n+1}=\{e\}$, es decir $\tau=\sigma$
 para todo conjugado $\tau$. Luego $\sigma\in Z(S_{n+1})=\{e\}$
 (para $n+1\ge 3$), contradicción. Por tanto, en este caso $N=\{e\}$.
 Así, los únicos subgrupos normales de $S_{n+1}$ son:
\[
  \{e\}\;(\text{orden }1),\quad
  A_{n+1}\;\bigl(\text{orden }\tfrac{(n+1)!}{2}\bigr),\quad
  S_{n+1}\;\bigl(\text{orden }(n+1)!\bigr).
\]
Como $n\ge 4$, se tiene $n!\ge 24$. Si $(n+1)!=n+1$, entonces
$n!=1$, imposible. Si $\tfrac{(n+1)!}{2}=n+1$, entonces
$\tfrac{n!}{2}=1$, es decir, $n!=2$, imposible. Por tanto ninguno
de estos órdenes coincide con $n+1$. Por tanto no existe subgrupo
normal de orden $n+1$ en $S_{n+1}$, y no puede existir tal homomorfismo
sobreyectivo. \hfill$\square$

\problem{2}{%
Sea $G$ un grupo conmutativo y sea $C$ el conjunto de clases conjugadas
de $G$, es decir, el conjunto de clases de equivalencia bajo la relación
\[
  g_1 \sim g_2 \quad \text{si existe } h \in G \text{ tal que } g_2 = h g_1 h^{-1}.
\]
Dada una clase conjugada $c \in C$, defina
\[
  s_c := \sum_{g \in c} g
\]
como un elemento en $\CC[G]$. Demuestre que $s_c$ es un elemento central,
es decir, que $s_c$ conmuta con todos los demás elementos del álgebra de grupo.
}

\answer
Basta demostrar que $h\,s_c = s_c\,h$ para todo $h\in G$,
pues los elementos $h\in G$ forman una base de $\CC[G]$ como $\CC$-espacio vectorial
y la igualdad se extiende por linealidad.

Equivalentemente, mostramos que $h\,s_c\,h^{-1}=s_c$.
Calculamos:
\[
  h\,s_c\,h^{-1}
  = h\Bigl(\sum_{g\in c}g\Bigr)h^{-1}
  = \sum_{g\in c}hgh^{-1}.
\]
Ahora bien, la conjugación por $h$ es una biyección $G\to G$
que preserva cada clase conjugada: si $g\in c$, entonces $hgh^{-1}\sim g$
(pues $hgh^{-1}$ es conjugado de $g$), así que $hgh^{-1}\in c$.
Por lo tanto la aplicación $g\mapsto hgh^{-1}$ es una biyección $c\to c$
(es inyectiva por ser restricción de una biyección, y $c$ es finito).

Al reindexar la suma con $g'=hgh^{-1}$, obtenemos
\[
  \sum_{g\in c}hgh^{-1}=\sum_{g'\in c}g'=s_c.
\]
Concluimos que $h\,s_c\,h^{-1}=s_c$, es decir $h\,s_c=s_c\,h$,
para todo $h\in G$. Esto muestra que $s_c\in Z(\CC[G])$. \hfill$\square$

\notehead{Observación} el enunciado dice ``grupo conmutativo'', pero la demostración
anterior vale para cualquier grupo finito $G$. Si $G$ es abeliano,
cada clase conjugada es un conjunto unitario $\{g\}$, $s_c=g$,
y el resultado se reduce a que $\CC[G]$ es conmutativo.

\problem{3}{%
Utilizando la terminología del ejercicio anterior, demuestre que
\[
  Z(\CC[G]) = \operatorname{span}_{\CC}\langle s_c : c \in C \rangle.
\]
Donde, si $A$ es un anillo, $Z(A)$ denota el centro de $A$, es decir,
el conjunto de elementos del álgebra que conmutan con todos los demás.
}

\answer

\subhead{Contención $\supseteq$}
$Z(\CC[G])$ es un $\CC$-subespacio de $\CC[G]$, y por el Problema~2 cada $s_c$
pertenece a $Z(\CC[G])$; luego
$\operatorname{span}_{\CC}\langle s_c : c\in C\rangle\subseteq Z(\CC[G])$.

\subhead{Contención $\subseteq$}
Sea $z=\sum_{g\in G}\alpha_g\,g\in Z(\CC[G])$.
Para todo $h\in G$ se tiene $hzh^{-1}=z$, es decir
\[
  \sum_{g\in G}\alpha_g\,hgh^{-1}
  =\sum_{g\in G}\alpha_g\,g.
\]
Sustituyendo $g'=hgh^{-1}$ en el lado izquierdo
(biyección $G\to G$, con $g=h^{-1}g'h$):
\[
  \sum_{g'\in G}\alpha_{h^{-1}g'h}\,g'
  =\sum_{g'\in G}\alpha_{g'}\,g'.
\]
Como los elementos de $G$ son base de $\CC[G]$,
la igualdad coeficiente a coeficiente da
\[
  \alpha_{h^{-1}g'h}=\alpha_{g'}\qquad
  \text{para todos }g'\in G,\;h\in G.
\]
Esto dice que la función $g\mapsto\alpha_g$ es constante
en cada clase conjugada.
Denotando por $\alpha_c$ el valor común para $g\in c$,
agrupamos la suma por clases conjugadas (que particionan $G$):
\[
  z=\sum_{c\in C}\;\sum_{g\in c}\alpha_g\,g
   =\sum_{c\in C}\alpha_c\sum_{g\in c}g
   =\sum_{c\in C}\alpha_c\,s_c
   \;\in\;\operatorname{span}_{\CC}\langle s_c:c\in C\rangle.
\]

De ambas contenciones,
$Z(\CC[G])=\operatorname{span}_{\CC}\langle s_c:c\in C\rangle$.
\hfill$\square$

\notehead{Observación} los $s_c$ son $\CC$-linealmente independientes,
pues cada $g\in G$ aparece en exactamente un $s_c$ (el de su clase conjugada);
por tanto, en una relación $\sum_c\lambda_c s_c=0$, el coeficiente de cualquier
$g\in c$ es $\lambda_c$, luego $\lambda_c=0$ para toda clase $c$.
Así, $\{s_c:c\in C\}$ es una $\CC$-base de $Z(\CC[G])$ y
$\dim_\CC Z(\CC[G])=|C|$.

\problem{4}{%
Escriba $\CC[C_4]$ como suma directa de irreps.
}

\answer
Sea $C_4=\langle r\mid r^4=e\rangle=\{e,r,r^2,r^3\}$ el grupo cíclico de orden~$4$.

Consideremos la representación regular
$\rho\colon C_4\to GL(\CC[C_4])\cong GL_4(\CC)$,
donde $g\in C_4$ actúa por multiplicación a izquierda sobre $\CC[C_4]$.
En la base ordenada $\mathcal B=(e,r,r^2,r^3)$, la matriz de $\rho(r)$ es
\[
  P=
  \begin{pmatrix}
    0&0&0&1\\
    1&0&0&0\\
    0&1&0&0\\
    0&0&1&0
  \end{pmatrix},
\]
ya que $r\cdot e=r$, $r\cdot r=r^2$, $r\cdot r^2=r^3$ y $r\cdot r^3=e$.

Como $r^4=e$, se tiene $P^4=I$. Por tanto el polinomio mínimo de $P$
divide a $x^4-1$, y sobre $\CC$ vale
\[
  x^4-1=(x-1)(x-i)(x+1)(x+i),
\]
que tiene cuatro raíces distintas. Luego $P$ es diagonalizable y sus autovalores
pertenecen a $\{1,i,-1,-i\}$.

Definimos los vectores
\[
  \begin{aligned}
    v_0&=\begin{pmatrix}1\\1\\1\\1\end{pmatrix},
    & v_1&=\begin{pmatrix}1\\-i\\-1\\i\end{pmatrix},\\[0.35em]
    v_2&=\begin{pmatrix}1\\-1\\1\\-1\end{pmatrix},
    & v_3&=\begin{pmatrix}1\\i\\-1\\-i\end{pmatrix}.
  \end{aligned}
\]
Se verifica directamente que
\[
  Pv_0=v_0,
  \qquad
  Pv_1=i\,v_1,
  \qquad
  Pv_2=-v_2,
  \qquad
  Pv_3=-i\,v_3.
\]
Así, los subespacios $\CC v_0$, $\CC v_1$, $\CC v_2$ y $\CC v_3$
son invariantes por $C_4$ y de dimensión~$1$; en particular, son irreducibles.

Si llamamos $\rho_k$ a la representación $1$-dimensional dada por
$\rho_k(r)=i^k$, con $k=0,1,2,3$, entonces
$(\CC v_k,\rho|_{\CC v_k})\cong (\CC,\rho_k)$.

\subhead{Conclusión}
Como representación de $C_4$, la descomposición de la representación regular es
\[
  (\CC[C_4],\rho)
  \;\cong\; (\CC,\rho_0)\oplus(\CC,\rho_1)\oplus(\CC,\rho_2)\oplus(\CC,\rho_3).
\]
Equivalentemente,
$\CC[C_4]=\CC v_0\oplus \CC v_1\oplus \CC v_2\oplus \CC v_3$,
y en cada sumando $\CC v_k$ la acción de $r$ es multiplicación por $i^k$.
\hfill$\square$

\problem{5}{%
Escriba $\CC[S_3]$ como suma directa de irreps.
}

\begingroup
\small
\setlength{\parskip}{0.25em}
\setlength{\abovedisplayskip}{0.35em}
\setlength{\belowdisplayskip}{0.35em}
\setlength{\abovedisplayshortskip}{0.2em}
\setlength{\belowdisplayshortskip}{0.2em}
\setlength{\jot}{0.15em}

\subhead{Hecho (Corolario~4.3.15 de Steinberg~\cite[Cor.~4.3.15]{Steinberg})}
Si $\chi$ es el carácter de una representación compleja de un grupo finito $G$,
entonces la representación es irreducible si y solo si
\[
  \langle \chi,\chi\rangle
  :=\frac{1}{|G|}\sum_{g\in G}\chi(g)\overline{\chi(g)}=1.
\]
Como los caracteres son constantes en clases conjugadas, la suma puede agruparse
por clases conjugadas.

\answer
Sea $P_3$ la representación de permutación de $S_3$ sobre $\CC^3$,
donde cada $\sigma\in S_3$ actúa permutando la base canónica.
Consideremos los subespacios
\[
  L:=\CC(1,1,1),
  \qquad
  V:=\{(z_1,z_2,z_3)\in\CC^3:z_1+z_2+z_3=0\}.
\]
Ambos son $S_3$-invariantes: $L$ porque permutar coordenadas deja fijo al vector
$(1,1,1)$, y $V$ porque permutar coordenadas no cambia la suma $z_1+z_2+z_3$.
Además,
\[
  \CC^3=L\oplus V,
\]
porque $L\cap V=\{0\}$ y $\dim L+\dim V=1+2=3$.
Por tanto,
\[
  P_3\cong \mathbf 1\oplus V,
\]
donde $\mathbf 1$ es la representación trivial y $V$ tiene dimensión $2$.

Veamos que $V$ es irreducible. El carácter de $P_3$ se obtiene contando
cuántos vectores de la base canónica quedan fijos:
\[
  \chi_{P_3}(e)=3,
  \qquad
  \chi_{P_3}(\tau)=1\ \text{si }\tau\text{ es una transposición},
  \qquad
  \chi_{P_3}(c)=0\ \text{si }c\text{ es un }3\text{-ciclo}.
\]
Como $P_3\cong \mathbf 1\oplus V$, se sigue que
\[
  \chi_V(e)=2,
  \qquad
  \chi_V(\tau)=0,
  \qquad
  \chi_V(c)=-1.
\]
Entonces
\[
  \langle \chi_V,\chi_V\rangle
  =\frac{1}{|S_3|}\Bigl(1\cdot 2^2+3\cdot 0^2+2\cdot (-1)^2\Bigr)
  =\frac{1}{6}(4+0+2)=1.
\]
Por el criterio de irreducibilidad por caracteres, $V$ es irreducible.

Determinemos ahora todas las irreps de $S_3$.
Como $S_3$ tiene tres clases conjugadas
($\{e\}$, $\{(1\;2),(1\;3),(2\;3)\}$, $\{(1\;2\;3),(1\;3\;2)\}$),
hay exactamente tres irreps sobre $\CC$
\cite[Thm.~4.4.8]{Steinberg}.
Toda representación de dimensión~$1$, $\rho\colon S_3\to\CC^{\times}$,
se factoriza por el abelianizado $S_3/[S_3,S_3]=S_3/A_3\cong C_2$,
pues $\CC^{\times}$ es abeliano y por tanto $[S_3,S_3]\subseteq\ker\rho$.
Así, los únicos homomorfismos $C_2\to\CC^{\times}$ dan las irreps
$\mathbf 1$ y $\operatorname{sgn}$.
La ecuación de dimensiones~\cite[Cor.~4.4.5]{Steinberg}
$1^2+1^2+d_3^2=|S_3|=6$ fuerza $d_3=2$, y esa tercera irrep es $V$.

Por la descomposición de la representación
regular~\cite[Thm.~4.4.4]{Steinberg},
$\CC[G]\cong\bigoplus_i d_iW_i$, obtenemos
\[
  \CC[S_3]\cong \mathbf 1\oplus \operatorname{sgn}\oplus V\oplus V.
\]

\medskip\noindent\textbf{Subespacios explícitos en $\CC[S_3]$.}
Escribamos la base canónica de $\CC[S_3]$ como
$\{e_\sigma:\sigma\in S_3\}$, con acción $g\cdot e_\sigma=e_{g\sigma}$.
Los generadores son $s=(1\;2)$ y $r=(1\;2\;3)$; un subespacio de
$\CC[S_3]$ es $S_3$-invariante si y solo si es invariante bajo $s$ y $r$.
Sea $v=\sum_\sigma a_\sigma\,e_\sigma$.

\smallskip\noindent\emph{(a) Copia de\/ $\mathbf 1$}
(resolver $g\cdot v=v$ para todo $g$, es decir $s\cdot v=v$ y $r\cdot v=v$).

La acción de $s$ intercambia $e\leftrightarrow(1\;2)$,
$(1\;3)\leftrightarrow(1\;3\;2)$, $(2\;3)\leftrightarrow(1\;2\;3)$.
La de $r$ cicla $e\to(1\;2\;3)\to(1\;3\;2)\to e$ y
$(1\;2)\to(1\;3)\to(2\;3)\to(1\;2)$.
Las ecuaciones $s\cdot v=v$ y $r\cdot v=v$ fuerzan
$a_\sigma=a_\tau$ para todo par $\sigma,\tau$, dando
el subespacio de dimensión~$1$
\[
  \CC\cdot v_{\mathbf 1},
  \qquad
  v_{\mathbf 1}=\sum_{\sigma\in S_3}e_\sigma.
\]

\smallskip\noindent\emph{(b) Copia de\/ $\operatorname{sgn}$}
(resolver $g\cdot v=\operatorname{sgn}(g)\,v$,
i.e.\ $s\cdot v=-v$ y $r\cdot v=v$).

$s\cdot v=-v$ da $a_{(12)}=-a_e$, $a_{(132)}=-a_{(13)}$,
$a_{(123)}=-a_{(23)}$; $r\cdot v=v$ da
$a_e=a_{(123)}=a_{(132)}$ y $a_{(12)}=a_{(13)}=a_{(23)}$.
Combinando: $a_{(12)}=a_{(13)}=a_{(23)}=-a_e=-a_{(123)}=-a_{(132)}$.
El subespacio es $\CC\cdot v_{\operatorname{sgn}}$ con
\[
  v_{\operatorname{sgn}}
    = e_e+e_{(123)}+e_{(132)}-e_{(12)}-e_{(13)}-e_{(23)}.
\]

\smallskip\noindent\emph{(c) Dos copias de\/ $V$}
(componente estándar, $\dim=4$).

El complemento de $\operatorname{span}\{v_{\mathbf 1},
v_{\operatorname{sgn}}\}$ en $\CC[S_3]$ es
\[
  U=\bigl\{{\textstyle\sum_\sigma} a_\sigma\,e_\sigma :
    a_e+a_{(123)}+a_{(132)}=0,\;
    a_{(12)}+a_{(13)}+a_{(23)}=0\bigr\},
\]
de dimensión $4$.
Para descomponerlo, diagonalizamos $r$ en $U$.
Como $r^3=e$, los eigenvalores posibles son $1,\omega,\bar\omega$
con $\omega=e^{2\pi i/3}$.
El eigenvalor $1$ no aparece en $U$
(las dos ecuaciones de $U$ junto con $r\cdot v=v$ fuerzan $v=0$).
Los de eigenvalor $\omega$ satisfacen $a_{(123)}=\omega\,a_e$,
$a_{(132)}=\omega^2 a_e$, $a_{(13)}=\omega\,a_{(12)}$,
$a_{(23)}=\omega^2 a_{(12)}$
(las ecuaciones de $U$ se cumplen por $1+\omega+\omega^2=0$).
Esto da el $\omega$-eigenespacio de dimensión $2$:
\[
  \varphi_1=e_e+\omega^2 e_{(123)}+\omega\,e_{(132)},
  \qquad
  \varphi_2=e_{(12)}+\omega^2 e_{(13)}+\omega\,e_{(23)}.
\]
Conjugando coeficientes se obtiene el $\bar\omega$-eigenespacio
$\{\bar\varphi_1,\bar\varphi_2\}$.

Ahora calculamos la acción de $s=(1\;2)$. Recordando
$s\colon e\leftrightarrow(1\;2)$,
$(1\;3)\leftrightarrow(1\;3\;2)$,
$(2\;3)\leftrightarrow(1\;2\;3)$:
\begin{align*}
  s\cdot\varphi_1
    &= e_{(12)}+\omega^2 e_{(23)}+\omega\,e_{(13)}
    = \bar\varphi_2,\\
  s\cdot\varphi_2
    &= e_e+\omega^2 e_{(132)}+\omega\,e_{(123)}
    = \bar\varphi_1.
\end{align*}
Luego $W_1:=\operatorname{span}\{\varphi_1,\bar\varphi_2\}$ y
$W_2:=\operatorname{span}\{\varphi_2,\bar\varphi_1\}$ son
$S_3$-invariantes (cerrados bajo $r$ y $s$), con $U=W_1\oplus W_2$.
En la base $(\varphi_1,\bar\varphi_2)$ de $W_1$:
\[
  [r]_{W_1}=\begin{pmatrix}\omega&0\\0&\bar\omega\end{pmatrix},
  \qquad
  [s]_{W_1}=\begin{pmatrix}0&1\\1&0\end{pmatrix}.
\]
Las trazas son $\omega+\bar\omega=-1=\chi_V(r)$ y $0=\chi_V(s)$,
así que $W_1\cong V$; análogamente $W_2\cong V$.
(La descomposición $W_1\oplus W_2$ no es canónica: depende de cómo
emparejamos vectores entre los eigenespacios de $r$.)

\hfill$\square$
\endgroup

\problem{6}{%
Recuerde que $P_n$ es la representación $\CC^n$ de $S_n$ donde cada
elemento actúa permutando los elementos de la base canónica. Verifique
que $P_4$ se descompone como la suma directa de la irrep trivial y una
irrep de dimensión $3$.
}

\answer
Consideremos los subespacios
\[
  L:=\CC(1,1,1,1),
  \qquad
  V:=\{(z_1,z_2,z_3,z_4)\in\CC^4:z_1+z_2+z_3+z_4=0\}.
\]
Ambos son $S_4$-invariantes: $L$ porque permutar coordenadas deja fijo
al vector $(1,1,1,1)$, y $V$ porque permutar coordenadas no cambia la suma.
Además, $L\cap V=\{0\}$ y $\dim L+\dim V=1+3=4$, así que
$P_4\cong \mathbf 1\oplus V$,
donde $\mathbf 1$ denota la representación trivial (sobre $L$).

\subhead{Irreducibilidad de $V$}
Usamos el criterio de irreducibilidad por caracteres
(Corolario~4.3.15 de Steinberg~\cite[Cor.~4.3.15]{Steinberg}):
$V$ es irreducible si y solo si $\langle\chi_V,\chi_V\rangle=1$.

El carácter de $P_4$ es $\chi_{P_4}(\sigma)=|\operatorname{Fix}(\sigma)|$
(número de puntos fijos). Como $P_4\cong\mathbf 1\oplus V$ y $\chi_{\mathbf 1}\equiv 1$,
se tiene $\chi_V(\sigma)=|\operatorname{Fix}(\sigma)|-1$, y por tanto
\[
  \langle\chi_V,\chi_V\rangle
  =\frac{1}{4!}\sum_{\sigma\in S_4}(|\operatorname{Fix}(\sigma)|-1)^2
  =\frac{1}{4!}\Bigl(\sum_{\sigma}|\operatorname{Fix}(\sigma)|^2
    -2\sum_{\sigma}|\operatorname{Fix}(\sigma)|+4!\Bigr).
\]
Las dos sumas se calculan por conteo directo, sin enumerar clases conjugadas:

\notehead{Primera suma}
$\displaystyle\sum_{\sigma\in S_4}|\operatorname{Fix}(\sigma)|
=\#\{(i,\sigma):\ \sigma(i)=i\}$.
Para cada $i\in\{1,2,3,4\}$ hay $3!=6$ permutaciones que fijan $i$,
así que la suma es $4\cdot 6=24=4!$.

\notehead{Segunda suma}
$\displaystyle\sum_{\sigma\in S_4}|\operatorname{Fix}(\sigma)|^2
=\#\{(i,j,\sigma):\ \sigma(i)=i,\;\sigma(j)=j\}$.
Separamos los pares $(i,j)$ con $i=j$ y $i\ne j$:
\begin{itemize}[nosep,leftmargin=1.5em]
\item $i=j$: hay $4$ opciones para $i$ y $3!=6$ permutaciones que fijan $i$.
Contribución: $4\cdot 6=24$.
\item $i\ne j$: hay $4\cdot 3=12$ pares y $2!=2$ permutaciones que fijan ambos.
Contribución: $12\cdot 2=24$.
\end{itemize}
Total: $24+24=48=2\cdot 4!$.

Sustituyendo:
\[
  \langle\chi_V,\chi_V\rangle
  =\frac{48-2\cdot 24+24}{24}=\frac{24}{24}=1.
\]
Por el criterio de irreducibilidad, $V$ es irreducible de dimensión~$3$.

\subhead{Conclusión}
$P_4\cong\mathbf 1\oplus V$,
donde $\mathbf 1$ es la irrep trivial y $V$ es una irrep de dimensión $3$.
\hfill$\square$

\smallskip
{\footnotesize\textit{Observación:} el mismo argumento muestra que
$P_n\cong\mathbf 1\oplus V_n$ con $V_n:=\ker(\text{suma})\subset\CC^n$
irreducible de dimensión $n-1$ para todo $n\ge 2$
(la representación estándar de $S_n$).}

\begin{thebibliography}{9}

\bibitem{Steinberg}
Benjamin Steinberg,
\textit{Representation Theory of Finite Groups: An Introductory Approach},
Universitext, Springer, 2012.

\end{thebibliography}

\end{document}
